\documentclass[12pt]{article}
\usepackage{amsmath}
\usepackage{amssymb}
\usepackage{geometry}
\usepackage{amsthm}
\usepackage{bm}
\geometry{a4paper, margin=1in}

\title{Homework Chapter 4: Linearly Independent Sets; Bases}
\author{Tashfeen Omran}
\date{September 2026}

\begin{document}

\maketitle
\tableofcontents
\newpage

\section{Part 1: Comprehensive Introduction, Context, and Prerequisites}

\subsection{Core Concepts}
In the study of vector spaces, we often want to find the smallest, most efficient set of vectors that can generate the entire space. This brings us to the fundamental concepts of \textbf{Linear Independence} and \textbf{Bases}. 

A set of vectors $\{\mathbf{v}_1, \dots, \mathbf{v}_p\}$ in a vector space $V$ is said to be \textbf{linearly independent} if the vector equation
\[ c_1\mathbf{v}_1 + c_2\mathbf{v}_2 + \dots + c_p\mathbf{v}_p = \mathbf{0} \]
has \textit{only} the trivial solution ($c_1 = 0, c_2 = 0, \dots, c_p = 0$). If there is a nontrivial solution (where at least one weight $c_i$ is not zero), the set is \textbf{linearly dependent}.

A \textbf{basis} for a vector space $V$ (or a subspace $H$) is a set of vectors $\mathcal{B}$ that satisfies two strict conditions:
1. $\mathcal{B}$ is a linearly independent set.
2. The subspace spanned by $\mathcal{B}$ coincides with $V$ (i.e., $V = \text{Span } \mathcal{B}$).

\subsection{Intuition and Derivation}
Think of a vector space as a vast 3D city. A "spanning set" is a set of directions that allows you to reach any point in the city. However, a spanning set might contain redundant directions (e.g., "Go 1 mile North" and "Go 2 miles North"). 
Linear independence ensures there is no redundancy. No vector in the set can be built using a combination of the others. Therefore, a \textbf{basis} is the "DNA" or the minimal building blocks of the vector space. It is an "efficient" spanning set containing no unnecessary vectors.

\subsection{Historical Context and Motivation}
The formalization of vector spaces and bases occurred in the late 19th and early 20th centuries, heavily driven by mathematicians like Hermann Grassmann and Giuseppe Peano. They were motivated by the need to generalize Euclidean geometry to higher dimensions and to systematically solve massive systems of linear equations. By defining a "basis," mathematicians could map complex, abstract spaces (like spaces of polynomials or continuous functions) into simple coordinate systems, making impossible calculations trivial.

\subsection{Key Formulas and Theorems}
\begin{itemize}
    \item \textbf{Linear Independence Equation:} $c_1\mathbf{v}_1 + \dots + c_p\mathbf{v}_p = \mathbf{0}$.
    \item \textbf{Spanning Set Theorem:} If $S = \{\mathbf{v}_1, \dots, \mathbf{v}_p\}$ spans $V$, and one vector $\mathbf{v}_k$ is a linear combination of the others, then removing $\mathbf{v}_k$ leaves a set that still spans $V$.
    \item \textbf{Basis for Col $A$:} The \textit{pivot columns} of a matrix $A$ form a basis for the column space of $A$.
    \item \textbf{Basis for Nul $A$:} The vectors attached to the free variables in the parametric vector form of the solution to $A\mathbf{x} = \mathbf{0}$ form a basis for the null space of $A$.
\end{itemize}

\subsection{Prerequisites}
To master this chapter, you must be perfectly fluent in:
1. \textbf{Row Reduction (Gaussian Elimination):} Reducing a matrix to its Reduced Row Echelon Form (RREF).
2. \textbf{Solving Homogeneous Systems:} Finding the parametric vector form of $A\mathbf{x} = \mathbf{0}$.
3. \textbf{Spanning Sets:} Understanding what it means for a set of vectors to span a space.

\newpage
\section{Part 2: Detailed Homework Solutions}

Based on the provided text, a strategic selection of problems has been chosen to cover every distinct problem archetype (Checking bases, Nul $A$, Col $A$, Row $A$, Subspaces, True/False, and Polynomial spaces).

\subsection{Problem 3: Checking for a Basis in $\mathbb{R}^3$}
\textbf{Question:} Determine if the set of vectors is a basis for $\mathbb{R}^3$. Justify your answer.
\[ \mathbf{v}_1 = \begin{bmatrix} 1 \\ 0 \\ -2 \end{bmatrix}, \quad \mathbf{v}_2 = \begin{bmatrix} 3 \\ 2 \\ -4 \end{bmatrix}, \quad \mathbf{v}_3 = \begin{bmatrix} -3 \\ -5 \\ 1 \end{bmatrix} \]

\textbf{Solution:}
To be a basis for $\mathbb{R}^3$, the three vectors must be linearly independent and span $\mathbb{R}^3$. We place them into a matrix $A$ and row reduce to find the number of pivots.
\[ A = \begin{bmatrix} 1 & 3 & -3 \\ 0 & 2 & -5 \\ -2 & -4 & 1 \end{bmatrix} \]
Perform row operations to reach echelon form:
\begin{align*}
R_3 &\leftarrow R_3 + 2R_1 \\
A &\sim \begin{bmatrix} 1 & 3 & -3 \\ 0 & 2 & -5 \\ 0 & 2 & -5 \end{bmatrix} \\
R_3 &\leftarrow R_3 - R_2 \\
A &\sim \begin{bmatrix} 1 & 3 & -3 \\ 0 & 2 & -5 \\ 0 & 0 & 0 \end{bmatrix}
\end{align*}
The matrix has only 2 pivot positions (in columns 1 and 2). Because there is not a pivot in every column, the columns are linearly dependent. Because there is not a pivot in every row, they do not span $\mathbb{R}^3$. 
\textbf{Final Answer:} The set is \textbf{not} a basis for $\mathbb{R}^3$.

\subsection{Problem 5: Checking for a Basis in $\mathbb{R}^3$}
\textbf{Question:} Determine if the set of vectors is a basis for $\mathbb{R}^3$.
\[ \begin{bmatrix} 1 \\ -3 \\ 0 \end{bmatrix}, \begin{bmatrix} -2 \\ 9 \\ 0 \end{bmatrix}, \begin{bmatrix} 0 \\ 0 \\ 0 \end{bmatrix}, \begin{bmatrix} 0 \\ -3 \\ 5 \end{bmatrix} \]

\textbf{Solution:}
We have a set of 4 vectors in $\mathbb{R}^3$. By Theorem 8 from Chapter 1, any set containing more vectors than there are entries in each vector ($p > n$, here $4 > 3$) is automatically linearly dependent. Furthermore, the set contains the zero vector, which guarantees linear dependence. 
Since the set is linearly dependent, it cannot be a basis. 
(Note: If we row reduce the matrix formed by these vectors, we would find 3 pivots, meaning they \textit{do} span $\mathbb{R}^3$, but they fail the independence condition).
\textbf{Final Answer:} The set is \textbf{not} a basis for $\mathbb{R}^3$ because it is linearly dependent.

\subsection{Problem 9: Finding a Basis for Nul $A$}
\textbf{Question:} Find a basis for the null space of the matrix:
\[ A = \begin{bmatrix} 1 & 0 & -3 & 2 \\ 0 & 1 & -5 & 4 \\ 3 & -2 & 1 & -2 \end{bmatrix} \]

\textbf{Solution:}
The null space of $A$ is the set of all solutions to $A\mathbf{x} = \mathbf{0}$. We augment $A$ with the zero vector and row reduce:
\begin{align*}
\begin{bmatrix} 1 & 0 & -3 & 2 & 0 \\ 0 & 1 & -5 & 4 & 0 \\ 3 & -2 & 1 & -2 & 0 \end{bmatrix}
\end{align*}
\begin{align*}
R_3 &\leftarrow R_3 - 3R_1 \\
&\sim \begin{bmatrix} 1 & 0 & -3 & 2 & 0 \\ 0 & 1 & -5 & 4 & 0 \\ 0 & -2 & 10 & -8 & 0 \end{bmatrix} \\
R_3 &\leftarrow R_3 + 2R_2 \\
&\sim \begin{bmatrix} 1 & 0 & -3 & 2 & 0 \\ 0 & 1 & -5 & 4 & 0 \\ 0 & 0 & 0 & 0 & 0 \end{bmatrix}
\end{align*}
The matrix is in RREF. The pivot variables are $x_1$ and $x_2$. The free variables are $x_3$ and $x_4$.
Writing the system of equations:
\begin{align*}
x_1 - 3x_3 + 2x_4 &= 0 \implies x_1 = 3x_3 - 2x_4 \\
x_2 - 5x_3 + 4x_4 &= 0 \implies x_2 = 5x_3 - 4x_4 \\
x_3 &= x_3 \quad \text{(free)} \\
x_4 &= x_4 \quad \text{(free)}
\end{align*}
Expressing the solution in parametric vector form:
\[ \mathbf{x} = \begin{bmatrix} x_1 \\ x_2 \\ x_3 \\ x_4 \end{bmatrix} = x_3 \begin{bmatrix} 3 \\ 5 \\ 1 \\ 0 \end{bmatrix} + x_4 \begin{bmatrix} -2 \\ -4 \\ 0 \\ 1 \end{bmatrix} \]
The vectors that multiply the free variables form the basis for Nul $A$.
\textbf{Final Answer:} Basis for Nul $A$ is $\left\{ \begin{bmatrix} 3 \\ 5 \\ 1 \\ 0 \end{bmatrix}, \begin{bmatrix} -2 \\ -4 \\ 0 \\ 1 \end{bmatrix} \right\}$.

\subsection{Problem 13: Bases for Nul $A$, Col $A$, and Row $A$}
\textbf{Question:} Assume $A$ is row equivalent to $B$. Find bases for Nul $A$, Col $A$, and Row $A$.
\[ A = \begin{bmatrix} -2 & 4 & -2 & -4 \\ 2 & -6 & -3 & 1 \\ -3 & 8 & 2 & -3 \end{bmatrix}, \quad B = \begin{bmatrix} 1 & 0 & 6 & 5 \\ 0 & 2 & 5 & 3 \\ 0 & 0 & 0 & 0 \end{bmatrix} \]

\textbf{Solution:}
\textbf{1. Basis for Nul $A$:}
Nul $A$ is the same as Nul $B$. We solve $B\mathbf{x} = \mathbf{0}$.
\begin{align*}
x_1 + 6x_3 + 5x_4 &= 0 \implies x_1 = -6x_3 - 5x_4 \\
2x_2 + 5x_3 + 3x_4 &= 0 \implies x_2 = -\frac{5}{2}x_3 - \frac{3}{2}x_4
\end{align*}
$x_3$ and $x_4$ are free.
\[ \mathbf{x} = x_3 \begin{bmatrix} -6 \\ -5/2 \\ 1 \\ 0 \end{bmatrix} + x_4 \begin{bmatrix} -5 \\ -3/2 \\ 0 \\ 1 \end{bmatrix} \]
To avoid fractions, we can scale these vectors by 2 (since any scalar multiple of a basis vector is still a valid basis vector for the null space).
Basis for Nul $A$: $\left\{ \begin{bmatrix} -12 \\ -5 \\ 2 \\ 0 \end{bmatrix}, \begin{bmatrix} -10 \\ -3 \\ 0 \\ 2 \end{bmatrix} \right\}$.

\textbf{2. Basis for Col $A$:}
The pivot columns of $B$ are columns 1 and 2. Therefore, the pivot columns of the \textit{original} matrix $A$ form a basis for Col $A$.
Basis for Col $A$: $\left\{ \begin{bmatrix} -2 \\ 2 \\ -3 \end{bmatrix}, \begin{bmatrix} 4 \\ -6 \\ 8 \end{bmatrix} \right\}$.

\textbf{3. Basis for Row $A$:}
Row operations do not change the row space. The nonzero rows of the echelon form $B$ form a basis for Row $A$.
Basis for Row $A$: $\left\{ (1, 0, 6, 5), (0, 2, 5, 3) \right\}$.

\subsection{Problem 15: Basis for a Spanned Space}
\textbf{Question:} Find a basis for the space spanned by the given vectors:
\[ \mathbf{v}_1 = \begin{bmatrix} 1 \\ 0 \\ -3 \\ 2 \end{bmatrix}, \mathbf{v}_2 = \begin{bmatrix} 0 \\ 1 \\ 2 \\ -3 \end{bmatrix}, \mathbf{v}_3 = \begin{bmatrix} -3 \\ -4 \\ 1 \\ 6 \end{bmatrix}, \mathbf{v}_4 = \begin{bmatrix} 1 \\ -3 \\ -8 \\ 7 \end{bmatrix}, \mathbf{v}_5 = \begin{bmatrix} 2 \\ 1 \\ -6 \\ 9 \end{bmatrix} \]

\textbf{Solution:}
The space spanned by these vectors is exactly the column space of the matrix $A = [\mathbf{v}_1 \dots \mathbf{v}_5]$. We must row reduce $A$ to find its pivot columns.
\[ A = \begin{bmatrix} 1 & 0 & -3 & 1 & 2 \\ 0 & 1 & -4 & -3 & 1 \\ -3 & 2 & 1 & -8 & -6 \\ 2 & -3 & 6 & 7 & 9 \end{bmatrix} \]
\begin{align*}
R_3 &\leftarrow R_3 + 3R_1 \\
R_4 &\leftarrow R_4 - 2R_1 \\
&\sim \begin{bmatrix} 1 & 0 & -3 & 1 & 2 \\ 0 & 1 & -4 & -3 & 1 \\ 0 & 2 & -8 & -5 & 0 \\ 0 & -3 & 12 & 5 & 5 \end{bmatrix} \\
R_3 &\leftarrow R_3 - 2R_2 \\
R_4 &\leftarrow R_4 + 3R_2 \\
&\sim \begin{bmatrix} 1 & 0 & -3 & 1 & 2 \\ 0 & 1 & -4 & -3 & 1 \\ 0 & 0 & 0 & 1 & -2 \\ 0 & 0 & 0 & -4 & 8 \end{bmatrix} \\
R_4 &\leftarrow R_4 + 4R_3 \\
&\sim \begin{bmatrix} 1 & 0 & -3 & 1 & 2 \\ 0 & 1 & -4 & -3 & 1 \\ 0 & 0 & 0 & 1 & -2 \\ 0 & 0 & 0 & 0 & 0 \end{bmatrix}
\end{align*}
The pivots are in columns 1, 2, and 4. By the Spanning Set Theorem, the corresponding columns of the original matrix form a basis.
\textbf{Final Answer:} Basis is $\left\{ \begin{bmatrix} 1 \\ 0 \\ -3 \\ 2 \end{bmatrix}, \begin{bmatrix} 0 \\ 1 \\ 2 \\ -3 \end{bmatrix}, \begin{bmatrix} 1 \\ -3 \\ -8 \\ 7 \end{bmatrix} \right\}$.

\subsection{Problem 21 \& 22: True/False Conceptual Analysis}
\textbf{Question 21:} (T/F) A single vector by itself is linearly dependent.
\textbf{Solution:} \textbf{False}. A set containing a single vector $\mathbf{v}$ is linearly independent if and only if $\mathbf{v} \neq \mathbf{0}$. The equation $c_1\mathbf{v} = \mathbf{0}$ has only the trivial solution $c_1 = 0$ as long as $\mathbf{v}$ is not the zero vector.

\textbf{Question 22:} (T/F) A linearly independent set in a subspace $H$ is a basis for $H$.
\textbf{Solution:} \textbf{False}. To be a basis for $H$, the set must be linearly independent \textit{and} it must span $H$. A linearly independent set might be too small to span the entire subspace.

\subsection{Problem 43: Basis in Polynomial Space $\mathbb{P}_3$}
\textbf{Question:} Consider the polynomials $\mathbf{p}_1(t) = 1 + t^2$ and $\mathbf{p}_2(t) = 1 - t^2$. Is $\{\mathbf{p}_1, \mathbf{p}_2\}$ a linearly independent set in $\mathbb{P}_3$? Why or why not?

\textbf{Solution:}
We set up the linear independence equation for polynomials:
\[ c_1\mathbf{p}_1(t) + c_2\mathbf{p}_2(t) = 0 \quad \text{for all } t \]
\[ c_1(1 + t^2) + c_2(1 - t^2) = 0 \]
Group the terms by powers of $t$:
\[ (c_1 + c_2) + (c_1 - c_2)t^2 = 0 \]
For a polynomial to be the zero polynomial, all of its coefficients must be zero. This gives a system of equations:
\begin{align*}
c_1 + c_2 &= 0 \\
c_1 - c_2 &= 0
\end{align*}
Adding the two equations yields $2c_1 = 0 \implies c_1 = 0$. Substituting back yields $c_2 = 0$.
Since the only solution is the trivial solution, the polynomials are linearly independent.
\textbf{Final Answer:} \textbf{Yes}, they are linearly independent because neither is a scalar multiple of the other, and $c_1\mathbf{p}_1 + c_2\mathbf{p}_2 = 0$ implies $c_1=c_2=0$.

\subsection{Problem 44: Finding a Basis for a Polynomial Span}
\textbf{Question:} Consider $\mathbf{p}_1(t) = 1+t$, $\mathbf{p}_2(t) = 1-t$, and $\mathbf{p}_3(t) = 2$. Find a basis for Span$\{\mathbf{p}_1, \mathbf{p}_2, \mathbf{p}_3\}$.

\textbf{Solution:}
We need to find an efficient spanning set. Let's check for linear dependence. Notice that:
\[ \mathbf{p}_1(t) + \mathbf{p}_2(t) = (1+t) + (1-t) = 2 = \mathbf{p}_3(t) \]
Because $\mathbf{p}_3$ is a linear combination of $\mathbf{p}_1$ and $\mathbf{p}_2$, the Spanning Set Theorem allows us to discard $\mathbf{p}_3$ without changing the span.
Now we check if $\{\mathbf{p}_1, \mathbf{p}_2\}$ is linearly independent.
\[ c_1(1+t) + c_2(1-t) = 0 \implies (c_1+c_2) + (c_1-c_2)t = 0 \]
This implies $c_1+c_2=0$ and $c_1-c_2=0$, which means $c_1=c_2=0$. They are linearly independent.
\textbf{Final Answer:} A basis for the span is $\{\mathbf{p}_1, \mathbf{p}_2\}$.

\newpage
\section{Part 3: In-Depth Analysis of Problems and Techniques}

\subsection{A) Problem Types and General Approach}
\begin{itemize}
    \item \textbf{Type 1: Verifying a Basis for $\mathbb{R}^n$ (Problems 3, 5):} 
    \textit{Strategy:} Place the vectors as columns in a matrix $A$. Row reduce to echelon form. If there is a pivot in every column, they are independent. If there is a pivot in every row, they span $\mathbb{R}^n$. If both are true, it is a basis.
    \item \textbf{Type 2: Finding a Basis for Nul $A$ (Problems 9, 13):}
    \textit{Strategy:} Augment the matrix with zeros and find the RREF. Write out the equations, express pivot variables in terms of free variables, and factor out the free variables to get the basis vectors.
    \item \textbf{Type 3: Finding a Basis for Col $A$ (Problems 13, 15):}
    \textit{Strategy:} Row reduce $A$ to find which columns contain pivots. \textbf{Crucial step:} Go back to the \textit{original} matrix $A$ and select those specific columns.
    \item \textbf{Type 4: Polynomial Spaces (Problems 43, 44):}
    \textit{Strategy:} Treat polynomials like vectors by grouping like terms (constants, $t$, $t^2$). Set the resulting coefficients to zero to build a standard linear system.
\end{itemize}

\subsection{B) Key Algebraic and Calculus Manipulations}
\begin{itemize}
    \item \textbf{The "Original Column" Rule:} In Problem 13 and 15, row operations change the column space. Therefore, the RREF matrix only tells you \textit{where} the pivots are. You must extract the basis vectors from the original matrix.
    \item \textbf{Scaling Null Space Vectors:} In Problem 13, the null space basis vectors contained fractions ($-5/2, -3/2$). Because a subspace is closed under scalar multiplication, we multiplied the basis vectors by 2 to clear the fractions, making the basis cleaner to work with.
    \item \textbf{Inspection for Dependence:} In Problem 44, rather than setting up a full coordinate mapping matrix, we used algebraic inspection to notice $\mathbf{p}_1 + \mathbf{p}_2 = \mathbf{p}_3$. This immediately allowed us to apply the Spanning Set Theorem and discard $\mathbf{p}_3$.
\end{itemize}

\newpage
\section{Part 4: "Cheatsheet" and Tips for Success}

\subsection{Summary of Important Rules}
\begin{itemize}
    \item \textbf{Basis Definition:} $\mathcal{B}$ is a basis if it is linearly independent AND spans the space.
    \item \textbf{Nul $A$ Basis:} Solve $A\mathbf{x} = \mathbf{0}$. The vectors attached to the free variables form the basis.
    \item \textbf{Col $A$ Basis:} The pivot columns of the \textbf{original} matrix $A$.
    \item \textbf{Row $A$ Basis:} The non-zero rows of the \textbf{echelon form} matrix.
    \item \textbf{Dimension Check:} In $\mathbb{R}^n$, any set of $p$ vectors where $p > n$ is automatically linearly dependent (cannot be a basis).
    \item \textbf{Zero Vector Rule:} Any set containing the zero vector $\mathbf{0}$ is automatically linearly dependent.
\end{itemize}

\subsection{Common Pitfalls and Traps}
\begin{itemize}
    \item \textbf{The RREF Column Trap:} The most common mistake in linear algebra is using the columns of the RREF matrix to form the basis for Col $A$. Row operations preserve row space and null space, but they \textit{destroy} the original column space. Always go back to the original matrix!
    \item \textbf{Confusing "Spans" with "Basis":} A set can span a space without being a basis (if it has redundant vectors). A set can be independent without being a basis (if it doesn't have enough vectors to span the space).
    \item \textbf{Evaluating Polynomials at a Point:} When checking if functions are independent, you cannot just plug in $t=1$ and see if it equals zero. The equation $c_1\mathbf{p}_1(t) + c_2\mathbf{p}_2(t) = 0$ must hold for \textit{all} $t$.
\end{itemize}

\newpage
\section{Part 5: Conceptual Synthesis and The "Big Picture"}

\subsection{A) Thematic Connections}
The core theme of this chapter is \textbf{Dimensionality Reduction and Efficiency}. In mathematics, we frequently deal with massive amounts of data or infinitely large spaces. A basis allows us to compress an infinite space into a finite set of generators. 
This is conceptually identical to prime factorization in number theory. Just as every integer can be uniquely broken down into a product of prime numbers, every vector in a vector space can be uniquely broken down into a linear combination of basis vectors. The basis vectors are the "prime numbers" of linear algebra.

\subsection{B) Forward and Backward Links}
\textbf{Backward Links:} This topic is the logical culmination of Chapter 1 (Linear Combinations and Span) and Chapter 2 (Matrix Inverses). We are finally giving a formal name ("Basis") to the minimal spanning sets we've been intuitively building.

\textbf{Forward Links:} Understanding bases is the absolute bedrock for the rest of linear algebra. 
1. \textbf{Dimension and Rank:} The number of vectors in a basis defines the "dimension" of the space.
2. \textbf{Coordinate Systems:} In Section 4.4, we will use a basis to create custom coordinate systems, allowing us to view complex spaces as if they were just standard $\mathbb{R}^n$.
3. \textbf{Eigenvectors and Diagonalization:} Later, we will search for a specific basis made entirely of eigenvectors, which makes matrix multiplication as simple as scalar multiplication.

\newpage
\section{Part 6: Abstract Real-World Modeling (AI, ML, \& Quant Finance)}

\subsection{A) Concrete Scenarios}
1. \textbf{Principal Component Analysis (PCA) in Machine Learning:} In ML, datasets often have thousands of features (dimensions), many of which are highly correlated (linearly dependent). PCA finds a new orthogonal \textit{basis} for the data that maximizes variance. By projecting the data onto the subspace spanned by the first few basis vectors (Principal Components), we achieve massive dimensionality reduction while retaining the core information.
2. \textbf{Arbitrage Pricing Theory (APT) in Quant Finance:} In quantitative finance, the returns of thousands of assets are modeled as linear combinations of a small set of macroeconomic risk factors. These risk factors form a \textit{basis} for the portfolio space. If a portfolio can be constructed that is linearly independent of the risk basis but yields a risk-free return greater than the bank rate, an arbitrage opportunity exists.

\subsection{B) Model Problem Setup: PCA and Basis Projection}
Let us rigorously define the mathematical framework for dimensionality reduction using a basis in Statistical Learning Theory.

Let $\mathbf{X}$ be an $n \times p$ design matrix representing $n$ observations of $p$ random variables, centered such that the empirical mean of each column is zero. We assume $p$ is very large (e.g., $p = 10,000$ pixels in an image).
We define the empirical covariance matrix as:
\[ \Sigma = \frac{1}{n-1} \mathbf{X}^T \mathbf{X} \]
$\Sigma$ is a $p \times p$ symmetric, positive semi-definite matrix. By the Spectral Theorem, $\Sigma$ possesses a set of $p$ orthonormal eigenvectors $\{\mathbf{v}_1, \mathbf{v}_2, \dots, \mathbf{v}_p\}$ which form a \textbf{basis} for $\mathbb{R}^p$.

We sort these basis vectors such that their corresponding eigenvalues satisfy $\lambda_1 \geq \lambda_2 \geq \dots \geq \lambda_p \geq 0$. 
To reduce the dimensionality of our model from $p$ to $k$ (where $k \ll p$), we define a subspace $H = \text{Span}\{\mathbf{v}_1, \dots, \mathbf{v}_k\}$. The set $\mathcal{B}_k = \{\mathbf{v}_1, \dots, \mathbf{v}_k\}$ is a linearly independent set and forms a basis for $H$.

The optimization problem solved by PCA is finding the projection matrix $\mathbf{P}_k$ that minimizes the reconstruction error (loss function):
\[ \mathcal{L}(\mathbf{P}_k) = \frac{1}{n} \sum_{i=1}^n \left\| \mathbf{x}_i - \mathbf{P}_k \mathbf{x}_i \right\|^2 \]
where $\mathbf{P}_k = \mathbf{V}_k \mathbf{V}_k^T$ and $\mathbf{V}_k$ is the matrix whose columns are our basis vectors $\mathcal{B}_k$. By understanding how to extract a basis, AI models can compress a 10,000-dimensional image into a 50-dimensional vector, making real-time facial recognition computationally feasible.

\newpage
\section{Part 7: Common Variations and Untested Concepts}

\subsection{Extending a Linearly Independent Set to a Basis}
A common exam question that did not appear in the provided homework is \textbf{extending} a set to form a basis. 
\textbf{Theorem:} Any linearly independent set in a finite-dimensional vector space $V$ can be expanded to a basis for $V$ by adding appropriate vectors.

\textbf{Example Problem:} Let $H = \mathbb{R}^3$. You are given a linearly independent set $S = \{\mathbf{v}_1, \mathbf{v}_2\}$ where $\mathbf{v}_1 = \begin{bmatrix} 1 \\ 1 \\ 0 \end{bmatrix}$ and $\mathbf{v}_2 = \begin{bmatrix} 0 \\ 1 \\ 1 \end{bmatrix}$. Extend $S$ to a basis for $\mathbb{R}^3$.

\textbf{Solution:}
We need one more vector, $\mathbf{v}_3$, such that $\{\mathbf{v}_1, \mathbf{v}_2, \mathbf{v}_3\}$ is linearly independent. A standard trick is to append the standard basis vectors of $\mathbb{R}^3$ ($\mathbf{e}_1, \mathbf{e}_2, \mathbf{e}_3$) to our set and row reduce to find the pivot columns.
Form the matrix $A = [\mathbf{v}_1, \mathbf{v}_2, \mathbf{e}_1, \mathbf{e}_2, \mathbf{e}_3]$:
\[ A = \begin{bmatrix} 1 & 0 & 1 & 0 & 0 \\ 1 & 1 & 0 & 1 & 0 \\ 0 & 1 & 0 & 0 & 1 \end{bmatrix} \]
Row reduce $A$:
\begin{align*}
R_2 &\leftarrow R_2 - R_1 \\
&\sim \begin{bmatrix} 1 & 0 & 1 & 0 & 0 \\ 0 & 1 & -1 & 1 & 0 \\ 0 & 1 & 0 & 0 & 1 \end{bmatrix} \\
R_3 &\leftarrow R_3 - R_2 \\
&\sim \begin{bmatrix} 1 & 0 & 1 & 0 & 0 \\ 0 & 1 & -1 & 1 & 0 \\ 0 & 0 & 1 & -1 & 1 \end{bmatrix}
\end{align*}
The pivots are in columns 1, 2, and 3. 
Therefore, the first three columns of the original matrix form a basis for $\mathbb{R}^3$. We have successfully extended our set by adding $\mathbf{e}_1$.
\textbf{Extended Basis:} $\left\{ \begin{bmatrix} 1 \\ 1 \\ 0 \end{bmatrix}, \begin{bmatrix} 0 \\ 1 \\ 1 \end{bmatrix}, \begin{bmatrix} 1 \\ 0 \\ 0 \end{bmatrix} \right\}$.

\newpage
\section{Part 8: Advanced Diagnostic Testing: "Find the Flaw"}

Below are five problems with solutions that contain a subtle, critical flaw. Your task as the grader is to identify the error and provide the correct solution.

\subsection{Problem 1: Basis for Column Space}
\textbf{Question:} Find a basis for the column space of $A = \begin{bmatrix} 1 & 2 & 3 \\ 2 & 4 & 7 \\ 3 & 6 & 10 \end{bmatrix}$.
\textbf{Flawed Solution:}
Row reduce $A$:
\[ \begin{bmatrix} 1 & 2 & 3 \\ 2 & 4 & 7 \\ 3 & 6 & 10 \end{bmatrix} \sim \begin{bmatrix} 1 & 2 & 3 \\ 0 & 0 & 1 \\ 0 & 0 & 1 \end{bmatrix} \sim \begin{bmatrix} 1 & 2 & 0 \\ 0 & 0 & 1 \\ 0 & 0 & 0 \end{bmatrix} \]
The pivot columns are columns 1 and 3. Therefore, the basis for Col $A$ is $\left\{ \begin{bmatrix} 1 \\ 0 \\ 0 \end{bmatrix}, \begin{bmatrix} 0 \\ 1 \\ 0 \end{bmatrix} \right\}$.

\textbf{Correction:}
\textbf{Error:} The student used the columns of the RREF matrix instead of the original matrix $A$. Row operations change the column space.
\textbf{Correct Step:} The pivots are indeed in columns 1 and 3. We must select columns 1 and 3 from the \textit{original} matrix $A$.
\textbf{Correct Answer:} Basis for Col $A$ is $\left\{ \begin{bmatrix} 1 \\ 2 \\ 3 \end{bmatrix}, \begin{bmatrix} 3 \\ 7 \\ 10 \end{bmatrix} \right\}$.

\subsection{Problem 2: Linear Independence in $\mathbb{R}^4$}
\textbf{Question:} Determine if the set of vectors is linearly independent in $\mathbb{R}^4$: 
$S = \left\{ \begin{bmatrix} 1 \\ 0 \\ 1 \\ 2 \end{bmatrix}, \begin{bmatrix} 0 \\ 1 \\ -1 \\ 3 \end{bmatrix}, \begin{bmatrix} 2 \\ 2 \\ 0 \\ 10 \end{bmatrix} \right\}$.
\textbf{Flawed Solution:}
We have a set of 3 vectors in $\mathbb{R}^4$. According to the theorem on dimensions, if $n > p$ (where $n=4$ is the dimension and $p=3$ is the number of vectors), the set is automatically linearly dependent. Therefore, $S$ is linearly dependent.

\textbf{Correction:}
\textbf{Error:} The student reversed the theorem. A set is automatically dependent if $p > n$ (more vectors than entries). If $p \leq n$, they \textit{can} be independent, and we must check manually.
\textbf{Correct Step:} Set up $c_1\mathbf{v}_1 + c_2\mathbf{v}_2 + c_3\mathbf{v}_3 = \mathbf{0}$ and row reduce. Notice that $2\mathbf{v}_1 + 2\mathbf{v}_2 = \begin{bmatrix} 2 \\ 2 \\ 0 \\ 10 \end{bmatrix} = \mathbf{v}_3$. Since $\mathbf{v}_3$ is a linear combination of the others, the set is dependent.
\textbf{Correct Answer:} The set is linearly dependent (but not for the reason stated in the flawed solution).

\subsection{Problem 3: Basis for Null Space}
\textbf{Question:} Find a basis for Nul $A$ where $A = \begin{bmatrix} 1 & -2 & 0 \\ 0 & 0 & 1 \end{bmatrix}$.
\textbf{Flawed Solution:}
The matrix is already in RREF. The equations are $x_1 - 2x_2 = 0$ and $x_3 = 0$.
This means $x_1 = 2x_2$. If we set the free variable $x_2 = 0$, we get $x_1 = 0$ and $x_3 = 0$.
Since the only solution is the zero vector, Nul $A = \{\mathbf{0}\}$, which has no basis.

\textbf{Correction:}
\textbf{Error:} The student set the free variable to 0, which only finds the trivial solution. To find the basis, you must factor out the free variable, not eliminate it.
\textbf{Correct Step:} Write the parametric vector form: $\mathbf{x} = \begin{bmatrix} 2x_2 \\ x_2 \\ 0 \end{bmatrix} = x_2 \begin{bmatrix} 2 \\ 1 \\ 0 \end{bmatrix}$.
\textbf{Correct Answer:} The basis for Nul $A$ is $\left\{ \begin{bmatrix} 2 \\ 1 \\ 0 \end{bmatrix} \right\}$.

\subsection{Problem 4: Independence of Polynomials}
\textbf{Question:} Are the polynomials $\mathbf{p}_1(t) = 1$, $\mathbf{p}_2(t) = t$, and $\mathbf{p}_3(t) = t^2$ linearly independent?
\textbf{Flawed Solution:}
Let's test if $c_1(1) + c_2(t) + c_3(t^2) = 0$.
If we plug in $t = 1$, we get $c_1 + c_2 + c_3 = 0$.
If we choose $c_1 = 1, c_2 = 1, c_3 = -2$, the equation holds ($1 + 1 - 2 = 0$).
Since we found non-zero weights $c_i$ that make the equation zero, the polynomials are linearly dependent.

\textbf{Correction:}
\textbf{Error:} Linear independence of functions requires $c_1\mathbf{p}_1(t) + c_2\mathbf{p}_2(t) + c_3\mathbf{p}_3(t) = 0$ for \textbf{all} $t$, not just a single specific value of $t$. The weights $c_i$ must work universally.
\textbf{Correct Step:} $c_1 + c_2t + c_3t^2 = 0$ for all $t$ implies that the coefficients of the polynomial on the left must match the coefficients of the zero polynomial on the right. Thus $c_1=0, c_2=0, c_3=0$.
\textbf{Correct Answer:} The polynomials are linearly independent.

\subsection{Problem 5: Spanning Set Theorem}
\textbf{Question:} Let $S = \{\mathbf{v}_1, \mathbf{v}_2, \mathbf{v}_3\}$ span a subspace $H$. If $\mathbf{v}_3$ is NOT a linear combination of $\mathbf{v}_1$ and $\mathbf{v}_2$, can we remove $\mathbf{v}_3$ and still span $H$?
\textbf{Flawed Solution:}
Yes. By the Spanning Set Theorem, we can always remove a vector from a spanning set to shrink it down to a basis. Therefore, Span$\{\mathbf{v}_1, \mathbf{v}_2\} = H$.

\textbf{Correction:}
\textbf{Error:} The Spanning Set Theorem explicitly states you can only remove a vector if it \textbf{is} a linear combination of the remaining vectors. Removing a vector that is independent of the others will shrink the span.
\textbf{Correct Step:} Since $\mathbf{v}_3$ is not a linear combination of the others, removing it means the new set will no longer span the original subspace $H$.
\textbf{Correct Answer:} No, removing $\mathbf{v}_3$ will result in a set that does not span $H$.

\end{document}